\begin{table}[!htbp]
    \centering
    \caption{Zmanjševanje količine podatkov od izvirne zbirke \MAHNOB{} do končnih časovnih oken, uporabljenih v eksperimentalnem cevovodu.}
    \label{tab:mahnob-predobdelava}
    \small
    \begin{tabularx}{\textwidth}{@{}>{\raggedright\arraybackslash}p{0.28\textwidth} c c >{\raggedright\arraybackslash}X@{}}
        \toprule
        \textbf{Korak} & \textbf{Vrstice} & \textbf{Udeleženci} & \textbf{Opis} \\
        \midrule
        Izbira izvirne zbirke \MAHNOB{} & -- & 30 & Večmodalna zbirka; izvirni članek v analizi uporabi 27 udeležencev. \\
        Izbira podmnožice s samoocenami čustvenega stanja & 3\,797\,165 & 24 & Ohranimo le posnetke iz eksperimenta vzbujanja čustev, ki imajo podatke sledilnika pogleda in samoocene čustvenega stanja. \\
        Izključitev 'slabih' udeležencev & 3\,535\,842 & 22 & Izključena P9 ter P12. P15 izključen že prej. \\
        Izbira vrstic z veljavnimi samoocenami & 2\,995\,632 & 22 & Ohranimo le vrstice s podatki, uporabljenimi za učenje modelov samoocen čustvenega stanja. \\
        Priprava podatkov za učenje & 2\,814\,005 & 22 & Odstranjene so vrstice brez signalov pogleda ali zenic. \\
        Filter osamelcev \texttt{q01}-\texttt{q99} & 2\,639\,048 & 22 & Robustno filtriranje štirih osnovnih signalov pogleda in zenic. \\
        Tvorba uporabnih 10-sekundnih časovnih oken & 5\,158 & 22 & Vsako okno postane en učni primer oziroma graf. \\
        \bottomrule
    \end{tabularx}
\end{table}
